% !TeX program = xelatex
\documentclass[
  a4paper,fontsize=11pt,parskip=half,headings=small,
  numbers=noenddot,bibliography=totoc
]{scrartcl}

\usepackage[margin=16mm,includeheadfoot,headheight=14pt,headsep=6mm,footskip=9mm]{geometry}
\usepackage{fontspec}
\usepackage{mathtools}
\usepackage{unicode-math}
\setmainfont{TeX Gyre Pagella}
\setsansfont{TeX Gyre Heros}[Scale=MatchLowercase]
\setmonofont{Latin Modern Mono}[Scale=MatchLowercase]
\setmathfont{TeX Gyre Pagella Math}
\usepackage[protrusion=true,expansion=false]{microtype}
\usepackage{tabularray}
\UseTblrLibrary{booktabs}
\usepackage{booktabs}
\usepackage{enumitem}
\usepackage{xcolor}
\usepackage{scrlayer-scrpage}
\usepackage[authoryear,round]{natbib}
\usepackage{xurl}
\usepackage[
  unicode,colorlinks=true,linkcolor=blue!50!black,
  citecolor=blue!50!black,urlcolor=blue!50!black,
  pdftitle={Combining IQ Scores: Analysis and a Corrected Statistical Treatment},
  pdfsubject={Composite-score standardization, population calibration, and extreme percentiles},
  pdfauthor={Prepared with OpenAI Codex},
  pdfkeywords={IQ, composite scores, covariance, normal distribution, range restriction}
]{hyperref}
\usepackage{bookmark}
\usepackage[nameinlink,noabbrev]{cleveref}

\DeclareMathOperator{\Var}{Var}
\DeclareMathOperator{\Corr}{Corr}
\newcommand{\E}{\operatorname{E}}
\newcommand{\T}{\mathsf{T}}
\newcommand{\given}{\,\vert\,}
\setlength{\emergencystretch}{1.5em}
\setlist[itemize]{leftmargin=*,itemsep=2pt,topsep=3pt}
\setkomafont{disposition}{\normalfont\sffamily\bfseries}
\setkomafont{pageheadfoot}{\normalfont\sffamily\small}
\clearpairofpagestyles
\ihead{Combining IQ scores}
\ohead{Analysis and corrected article}
\cfoot*{\pagemark}
\pagestyle{scrheadings}
\setlength{\bibsep}{5pt}

\title{Combining IQ Scores}
\subtitle{Analysis and a corrected statistical treatment}
\author{}
\date{31 October 2026}

\begin{document}
\maketitle
\thispagestyle{plain.scrheadings}
\vspace{-1.1em}

\begin{abstract}
The covariance identities reproduced in Ronald K. Hoeflin's article are
correct. Their application to qualifying IQ scores requires additional
assumptions about the reference population, the composite's distribution,
and the meaning of the score. This report distinguishes those assumptions
from the algebra, corrects the numerical table, and supplies an independently
written replacement article. All displayed numerical examples have been
recalculated. A technical appendix illustrates why a standardized composite
and an estimate of latent ability need not coincide.
\end{abstract}

\section{Analysis of the original article}\label{sec:analysis}

\citet{hoeflin} reproduces George A. Ferguson's treatment of variances and
covariances and applies it to combinations of IQ scores. The basic identity
\begin{equation}
  \Var(X+Y)=\Var(X)+\Var(Y)+2\operatorname{Cov}(X,Y)
  \label{eq:variance-sum}
\end{equation}
is correct. It determines the variance of a sum without requiring normality.
The principal weaknesses arise when a standardized sum is assigned an extreme
population percentile and used as a qualifying score.

\begin{itemize}
  \item \textbf{Population calibration.} Correlations among the article's
  27 selected high scorers do not establish the correlations in the population
  used to standardize the tests. Selection can change the relationships, and
  the appropriate correction depends on how selection occurred
  \citep{sackett2000}.
  \item \textbf{Distributional assumptions.} Mean zero and variance one do
  not establish a normal distribution. A one-in-a-million upper-tail claim
  requires additional justification for the tail model.
  \item \textbf{Numerical correction.} The 100-test entry should be about
  3.69, rather than 3.71, standard deviations. The more accurate normal
  quantile also changes some final digits elsewhere; for example, the
  four-test result rounds to 3.98 rather than 3.97.
  \item \textbf{Score interpretation.} A standardized composite, a latent
  ability estimate, and a validated IQ score are distinct quantities.
  Their equivalence requires evidence about the battery and its intended use
  \citep[Standard~1.14]{standards2014}.
  \item \textbf{Threshold precision.} Rounded whole-number scores are useful
  illustrations, but are not exact qualifying cutoffs. The scale's standard
  deviation and the classification rule must be specified.
\end{itemize}

The replacement below preserves the valid calculation while making its
assumptions explicit. The article is newly written; it is not a transcription
or a reproduction of the original scans.

\section{Corrected article: a statistical method and its limits}
\label{sec:replacement}

Several cognitive test scores can be combined into a single score. The
calculation must account for the tests' correlations: correlated scores
cannot be treated as independent when calculating the composite's variance.
The resulting composite also needs an explicit interpretation. Rescaling it
to an IQ-style scale supplies units; evidence about the battery establishes
what those units mean.

Before calculating a composite, specify the tests, their weights, the
administration procedure, and the population against which scores will be
compared. Composite construction and interpretation require their own
justification; evidence supporting the individual tests does not
automatically establish the meaning of their combination
\citep[Standard~1.14]{standards2014}.

\subsection{Put the component scores on comparable scales}

For test \(i\), convert the observed score \(x_i\) to a standardized score:
\begin{equation}
  z_i=\frac{x_i-\mu_i}{\sigma_i},
  \label{eq:standardize}
\end{equation}
where \(\mu_i\) and \(\sigma_i\) are its mean and standard deviation in the
intended reference population. The norms must support comparison with that
population. Using different numbers labelled ``IQ'' does not establish
their comparability.

For example, a score of 160 on a scale with mean 100 and standard deviation
15 corresponds to \(z=4\). The same numerical score on a scale with standard
deviation 16 corresponds to \(z=3.75\). These scores should be standardized
before they are combined.

\subsection{Standardize the weighted composite}

Let \(R\) be the population correlation matrix of the standardized test
scores, and let \(w\) contain fixed weights. The weighted sum has variance
\begin{equation}
  \Var(w^\T Z)=w^\T Rw.
  \label{eq:weighted-variance}
\end{equation}
Consequently, provided this variance is positive, a standardized composite is
\begin{equation}
  \boxed{Z_C=\frac{w^\T z}{\sqrt{w^\T Rw}}.}
  \label{eq:weighted-composite}
\end{equation}
With the correct population parameters, this composite has mean zero and
variance one. This statement requires finite variances, but does not require
normality. When the parameters are estimated, their uncertainty affects the
calibration of the composite.

For an equally weighted battery of \(n\geq2\) tests, write \(\bar z\) for the
average standardized score and \(\bar r\) for the average correlation across
the distinct pairs of tests:
\begin{equation}
  \bar r=\frac{2}{n(n-1)}\sum_{i<j}r_{ij}.
  \label{eq:average-correlation}
\end{equation}
\begin{samepage}
Then
\begin{equation}
  \begin{aligned}
    \Var(\bar Z)
    &=\frac{n+2\sum_{i<j}r_{ij}}{n^2}\\
    &=\frac{1+(n-1)\bar r}{n}
     =\bar r+\frac{1-\bar r}{n}.
  \end{aligned}
  \label{eq:mean-variance}
\end{equation}
Therefore,
\begin{equation}
  \boxed{Z_C=\frac{\bar z}{\sqrt{\bar r+(1-\bar r)/n}}.}
  \label{eq:equal-composite}
\end{equation}
\end{samepage}
For a single test, \(Z_C=z_1\); no inter-test correlation is needed.

Identical pairwise correlations are unnecessary for
\cref{eq:mean-variance}; their average is sufficient for this particular
equal-weight composite. The full matrix remains relevant to other weighting
schemes and to understanding the battery.

An IQ-style reporting scale with mean 100 and standard deviation \(s\) can
be defined by
\begin{equation}
  I_C=100+sZ_C.
  \label{eq:reporting-scale}
\end{equation}
When all component scores already use mean 100 and the same standard
deviation \(s\), the formula becomes
\begin{equation}
  I_C=100+\frac{\bar I-100}{\sqrt{\bar r+(1-\bar r)/n}},
  \label{eq:iq-composite}
\end{equation}
where \(\bar I\) is their arithmetic average.

For example, three scores averaging 160, each on an SD-16 scale, produce
\(\bar z=3.75\). If their population average correlation is 0.60, then
\begin{equation}
  Z_C=\frac{3.75}{\sqrt{0.60+0.40/3}}\approx4.3791,
  \qquad I_C\approx170.06.
  \label{eq:worked-example}
\end{equation}
This is a score on the specified composite. Interpreting it as a measure of
general intelligence requires evidence about what the composite measures.

\subsection{Separate standardization from percentile estimation}

To assign a percentile, we also need the composite's distribution. Mean zero
and variance one do not make a distribution normal. Joint multivariate
normality of the component scores is a sufficient condition for a fixed
linear composite to be normal. Normal distributions for the individual
tests alone are insufficient.

If a standard-normal distribution is justified for the composite, its
percentile is \(100\Phi(Z_C)\), where \(\Phi\) is the standard-normal
cumulative distribution function. A one-in-a-million upper-tail probability
corresponds to
\begin{equation}
  q=\Phi^{-1}(0.999999)\approx4.7534243088.
  \label{eq:quantile}
\end{equation}
Thus the model-based 99.9999th-percentile threshold is
\begin{equation}
  \boxed{\bar z_{\min}=q\sqrt{\bar r+\frac{1-\bar r}{n}}.}
  \label{eq:threshold}
\end{equation}
The quantile can be reproduced in R with
\texttt{qnorm(1e-6, lower.tail = FALSE)}; the official documentation defines
the normal quantile function and its upper-tail option \citep{rnormal}.

For \(\bar r=0.60\), the thresholds are shown in \cref{tab:thresholds}.

\begin{table}[htbp]
\centering
\caption{Model-based 99.9999th-percentile thresholds, assuming a normal
composite and average pairwise correlation 0.60.}
\label{tab:thresholds}
\begin{tblr}{
  colspec={X[1,l] X[1.15,r] X[1.4,r]},
  row{1}={font=\bfseries},
  rowsep=4pt,
}
\toprule
Number of tests & Required average component \(z\)
  & Required average component score, mean 100 and SD 16 \\
\midrule
1     & 4.7534 & 176.05 \\
2     & 4.2516 & 168.03 \\
3     & 4.0706 & 165.13 \\
4     & 3.9770 & 163.63 \\
5     & 3.9198 & 162.72 \\
10    & 3.8027 & 160.84 \\
100   & 3.6942 & 159.11 \\
1,000 & 3.6832 & 158.93 \\
Limit as \(n\) increases & 3.6820 & 158.91 \\
\bottomrule
\end{tblr}
\end{table}

These are thresholds for the \emph{average component score}, under the
stated model. The target composite score remains approximately 176.05 on
the SD-16 scale in every row. Values are rounded for display; an operational
classification rule must specify rounding and use the appropriate
attainable scores.

For an SD-15 reporting scale, use \(100+15\bar z_{\min}\) instead. In
particular, the single-test threshold would be approximately 171.30. An IQ
number without its scale convention is ambiguous.

\subsection{Interpret the limiting threshold and correlation sensitivity}

The limiting average threshold follows from
\begin{equation}
  \lim_{n\to\infty}\bar z_{\min}=q\sqrt{\bar r},
  \label{eq:limit}
\end{equation}
if the average correlation remains fixed as tests are added. With
\(\bar r=0.60\), the limiting average component score on the SD-16 scale is
approximately 158.91. This is a feature of the assumed sequence of batteries
and its threshold; it is not a ceiling on possible composite scores. Adding
tests can also change the average correlation, so the fixed-correlation
limit should not be treated as a universal rule.

The choice of correlation matters substantially. For three tests, an
average correlation of 0.90 raises the required average SD-16 component
score to approximately 173.48, compared with 165.13 at 0.60. A smaller
assumed correlation gives a larger composite score for the same positive
average performance. The correlation therefore needs to be estimated for
the relevant population and administration procedure. A lower correlation
does not, by itself, establish better measurement; it can also reflect
different constructs or greater measurement noise.

\subsection{Calibrate the population and the scoring procedure}

Correlations in a selected high-scoring sample need not equal correlations
in the population used to standardize the tests. Selection can change score
variances and relationships. The appropriate treatment depends on the
selection mechanism and the information available about the unrestricted
population; there is no universal adjustment. \citet{sackett2000} explain
why range-restriction corrections depend on those details. A small selected
sample cannot directly establish a one-in-a-million population frequency.

The scoring procedure also matters. A fixed battery, the best few results
from many attempts, and a rule that keeps adding tests until someone
qualifies have different score distributions. Each procedure requires
calibration for the rule actually used. Repeated administrations can
introduce practice effects, as documented by \citet{hausknecht2007}.
Repeated or overlapping tests should not be counted as fresh independent
observations without accounting for their dependence and administration
history.

Finally, variability across people is different from uncertainty about one
person. The correlation matrix determines how the composite varies in the
reference population. It does not, by itself, provide an individual
confidence interval, separate general ability from shared methods, or
establish measurement precision at extreme scores. Those questions require
an appropriate measurement model and supporting data.

A defensible report should identify the battery, norms, weights, estimated
correlations, and uncertainty. Direct composite norms should be used when
they adequately support the intended interpretation. If an extreme
percentile comes from a distributional model beyond the available normative
evidence, report it as a model-based extrapolation and show how sensitive
it is to the assumptions. The testing standards require clear descriptions
of reference populations and norming procedures in Standards~5.8--5.9,
and justification for score interpretations in Standards~1.1--1.4
\citep{standards2014}.

\appendix
\section{Composite scores and latent ability: a mathematical illustration}
\label{sec:latent}

To see why these quantities differ, suppose, expressly as a hypothetical
model, that
\begin{equation}
  Z_i=\sqrt{\rho}\,G+\sqrt{1-\rho}\,\varepsilon_i,
  \qquad 0<\rho<1,
  \label{eq:factor-model}
\end{equation}
where \(G\) and all \(\varepsilon_i\) are independent standard-normal
variables. This model treats all shared variance as one common factor and
all residuals as mutually independent. Real batteries need not satisfy
those assumptions.

In this particular model, the equal-weight composite's squared correlation
with \(G\) is
\begin{equation}
  \mathcal R=\Corr(G,Z_C)^2=\frac{n\rho}{1+(n-1)\rho}.
  \label{eq:factor-reliability}
\end{equation}
The conditional distribution of the latent factor has
\begin{align}
  \E[G\given Z_1,\ldots,Z_n]&=\sqrt{\mathcal R}\,Z_C,
  \label{eq:latent-mean}\\
  \Var(G\given Z_1,\ldots,Z_n)&=\frac{1-\rho}{1+(n-1)\rho}.
  \label{eq:latent-variance}
\end{align}
These identities follow from Gaussian conditioning. In particular,
\(\operatorname{Cov}(G,\bar Z)=\sqrt{\rho}\) and
\(\Var(\bar Z)=\rho+(1-\rho)/n\); the residual contrasts contain no
additional information about \(G\) under this model.

For three tests with \(\rho=0.60\), a composite at \(Z_C=4.7534\) has
conditional mean \(G\) of approximately 4.30 and conditional standard
deviation approximately 0.43. On a hypothetical latent scale with mean 100
and SD 16, these become a mean of approximately 168.79 and a standard
deviation of 6.82. The composite's reporting-scale value is approximately
176.05. The two values answer different questions, even in this unusually
simple model.

The illustration is a mathematical counterexample to treating a
standardized composite as automatically identical to a latent-ability
estimate. It is not a proposed correction factor for actual IQ scores,
and inter-test correlation should not generally be equated with reliability
or common-factor variance.

\section{Numerical reproduction}

The accompanying script \texttt{scripts/verify\_calculations.py} recalculates
the normal quantile, checks the table against \texttt{data/thresholds.csv}
and the numerical entries in this source, and verifies the worked examples.
It uses the upper-tail probability directly through
\texttt{scipy.stats.norm.isf(1e-6)}. The mathematical results do not depend
on the choice of software.

\section{Final software verification and its limits}
\label{sec:verification}

The final checks were executed on \texttt{computer10} on 31 October 2026.
The following results refer to the source and evidence preserved at
repository commit \texttt{2636a42}. The
\href{https://github.com/AndersH3/composite_iq_checked_claimes/blob/2636a42ef45dc3579eb0dafe10707a0fbec48b10/verification/evidence/deal-lint-clean-computer10-20261031/REPORT.md}{Deal verification report}
and the
\href{https://github.com/AndersH3/composite_iq_checked_claimes/blob/2636a42ef45dc3579eb0dafe10707a0fbec48b10/verification/evidence/deal-lint-spark-computer10-20261031/REPORT.md}{SPARK integration report}
link the exact commands, tool versions, source hashes, checksums, and raw
results. Repository-wide Python counts below cover both analyses and
their shared tooling; they are not counts for this article alone.

\subsection{Deal contracts and Python checks}

The contract-coverage gate found 20 Python files, 142 named functions, and
416 Deal decorators, with zero coverage gaps. All 19 positive files,
including tests and the gate itself, passed static lint with zero
diagnostics, empty output, and exit status 0, both normally and under
\texttt{python -O}. The deliberately incorrect fixture was checked
separately: it produced exactly the expected \texttt{DEL012}
postcondition diagnostic and exit status 1. Its failure is a required
negative control, not an unresolved production finding.

All 14 earlier production diagnostics were resolved. Bounded quotient
helpers retain the original division operations, numerical tolerances
remain unchanged, and the fixes do not disable contracts or suppress
diagnostics. The runtime, property, and generated-test suite passed
110 tests with no failures or skips. All 27 recorded command checks had
their expected outcomes, including numerical scripts, Ada-output
validators, replay, optimized-Python execution, and evidence audits.

Deal Solver proved the two pure proof-acceptance gate postconditions;
the deliberately wrong symbolic contract was rejected. This focused
symbolic result does not prove the complete Python programs. A clean
linter result is also not a universal proof: input-dependent contracts
remain runtime checks, and finite tests cannot establish correctness for
every possible input. The recorded environment used Python 3.13.5,
Deal 4.24.6, and Deal Solver 0.1.2.

\subsection{SPARK results for the Hoeflin/Ferguson core}

The fresh five-prover portfolio completed all 47 native checks for this
analysis: 8 flow checks and 39 proof checks, with zero
unproved and zero justified checks. CVC5, Z3, Alt-Ergo, classic COLIBRI,
and Yices were requested together. This is a portfolio result; it does
not mean that each prover independently discharged every obligation.
Incomplete individual runs remain identified in the earlier evidence.
The deliberately incorrect inclusion-exclusion control was rejected
with an unproved functional contract.

The proof project selects only \texttt{hoeflin\_core.ads} and
\texttt{hoeflin\_core.adb}. Under their stated input preconditions, the
contracts establish exact integer variance bookkeeping, bounded
floating-point transforms, and nested event counts with their specified
increments. The core uses no project assumptions, imported proof axioms,
or suppressed checks. The
\href{https://github.com/AndersH3/composite_iq_checked_claimes/blob/2636a42ef45dc3579eb0dafe10707a0fbec48b10/verification/evidence/deal-lint-spark-computer10-20261031/hoeflin/portfolio/gnatprove.out}{native Hoeflin proof report}
records the totals. The toolchain manifest records the compiler banner
as GNAT 13.2.0 and the GNATprove banner as \texttt{0.0w}; these distinct
identifiers should not be conflated.

\subsection{Numerical, simulation, and model boundaries}

The Ada drivers and \texttt{normal\_support} are outside SPARK. Their
elementary functions, C \texttt{erfc}, and inverse-tail bisection are
checked numerically against SciPy, not certified by the core proof.
Likewise, a proof of a floating-point bound does not establish an
approximation-error bound. The proof-friendly evaluation orders for
thresholds and composites are algebraically equivalent to the report's
formulas, but are not claimed to give bit-identical floating-point
results.

The retained simulation uses 1,000,000 trials for each of eight test-count
cases, with seed \texttt{20261030} reset across cases. These common random
numbers make the cases dependent. Independent Python replay reproduced
the tail counts and checked the means and variances within the recorded
tolerances, normally and under optimized Python. This establishes
reproducibility and numerical agreement, not independence of the
pseudo-random draws. Counts beyond the one-in-a-million cutoff remain
descriptive; this sample size does not validate that extreme tail.

Neither Deal nor SPARK proves the Gaussian common-factor model, the
population covariance calibration, an identification of correlation with
reliability, or the empirical validity of a psychometric interpretation.
Those assumptions and the distinction between a standardized composite
and latent ability remain as discussed above.

From a full Git clone, install the pinned development dependencies and
run \texttt{make check} at the repository root. The Python auditor
\path{verification/verify_evidence.py} checks the retained hashes, native
totals, and negative-control result. Historical Python identities are
resolved through verified Git blobs; changed Ada sources cannot use that
fallback. An evidence audit verifies record integrity and source identity;
it does not rerun the provers or independently validate their SMT proofs.

\bibliographystyle{plainnat}
\bibliography{references}

\end{document}
